\section{Analizando el Algoritmo}

complejidad
	. explico la complejidad
	. la demuestro
	. pongo ejemplos

Para analizar el algoritmo propuesto anteriormente, se puede ocurrir la manera de poner un ejemplo para analizar el comportamiento:
Siendo monedas = [1, 3, 4, 9, 6, 2]
\begin{enumerate}
	\item Turno de Sophia, ella sólo ve el par [1, 2] y agarra la moneda de mayor valor (el 2). Le toca ahora a Mateo, par [1, 6] y Sophia elige por él la moneda de valor 1. Sophia de vuelta [3, 6] elige el 6. Mateo [3, 9], suma 3. Sophia [4, 9] elige la moneda de valor 9. A Mateo finalmente le toca la moneda de valor 4. No quedan más monedas para elegir.
	\item Entonces, si se suman los valores de todas las monedas seleccionadas para cada uno: Puntos de Sophia = 2 + 6 + 9 = 17. Puntos de Mateo = 1 + 3 + 4 = 8. 
	\item Sophia gana.
\end{enumerate}
	
	
	
\subsection{Complejidad}

Cuando hablamos de \textit{complejidad algorítmica},

\subsection{Demostración}
POR INDUCCION

DEMOSTRACION
	. hago la demostracion matematica del por qué es optimo. demostracion por induccion. que sea precisa y clara.
	. hago la demostracion del empate.
	. mostrar ejemplos (5 creo que esta bien, 3 bordes, 2 normales) y como funciona


\subsection{Experimentos}

	. pongo mediciones en tiempo de distintos casos (hago el grafico y pongo la foto.png)
		. tiempos del algoritmo. 
		. optimalidad del algoritmo. (que es optimalidad ¿?)
